\documentclass{beamer}
\beamertemplatenavigationsymbolsempty
%\usepackage{enumitem}
\mode<presentation>
{
  %\usetheme{Madrid}
  \usetheme{CambridgeUS}
  %\usecolortheme{seahorse}
  \setbeamercovered{transparent}
}
\usepackage{graphicx}
\usepackage[english]{babel}
\usepackage[latin1]{inputenc}
\usepackage{times}
\usepackage[T1]{fontenc} 
% Or whatever. Note that the encoding and the font should match. If T1
% does not look nice, try deleting the line with the fontenc.
\usepackage{amsmath}
\usepackage{color}
\newcommand{\linespace}{\vskip 0.25cm}
\usepackage{rotating}

\definecolor{MyForestGreen}{rgb}{0,0.7,0} 
\newcommand{\tableemph}[1]{{#1}}
\newcommand{\tablewin}[1]{\tableemph{#1}}
\newcommand{\tablemid}[1]{\tableemph{#1}}
\newcommand{\tablelose}[1]{\tableemph{#1}}

\definecolor{MyLightGray}{rgb}{0.6,0.6,0.6}
\newcommand{\tabletie}[1]{\color{MyLightGray} {#1}}

% The text in square brackets is the short version of your title and will be used in the
% header/footer depending on your theme.
\title[Visualizing GP Ancestries]{Visualizing Genetic Programming Ancestries}

% Sub-titles are optional - uncomment and edit the next line if you want one.
% \subtitle{Why does sub-tree crossover work?} 

% The text in square brackets is the short version of your name(s) and will be used in the
% header/footer depending on your theme.
\author[McPhee, Casale, et al]{Nic McPhee\inst{1}, Maggie Casale\inst{2}, Finzel, Helmuth, \& Spector}

% The text in square brackets is the short version of your institution and will be used in the
% header/footer depending on your theme.
\institute[]
{
	\inst{1} Division of Science and Mathematics \\
	University of Minnesota, Morris \\
	Morris, Minnesota, USA \and %
	\inst{2} Design Center Inc. \\
	St. Paul, Minnesota, USA}

% The text in square brackets is the short version of the date if you need that.
\date{21 July 2016 \\ VizGEC}

% Delete this, if you do not want the table of contents to pop up at
% the beginning of each subsection:
\AtBeginSection[]
{
  \begin{frame}<beamer>
    \frametitle{Outline}
    \tableofcontents[currentsection, hideothersubsections]
  \end{frame}
}

\begin{document}

\begin{frame}
  \titlepage
\end{frame}

\section*{The Big Picture}




\subsection*{Silico-paleontology}

\begin{frame}{The great waste}
\begin{columns}
% Column 1
\begin{column}{0.6\textwidth}
\begin{itemize}
	\item Genetic programming clearly \emph{works}
    \item But we rarely know \emph{why} or \emph{how}
    \item Can generate \emph{huge} amounts of data
    \item We typically throw it away -- \& paleontologists weep!
    \item Silico-paleontology can help us understand and improve our tools
\end{itemize}
\end{column}
% Column 2
\begin{column}{0.4\textwidth}
% \center \includegraphics[height=0.75\textheight]{Illustrations/Trash_dump.jpg} \\
\tiny \url{https://www.flickr.com/photos/thibaud_saintin/12361648674}
\end{column}
\end{columns}
\end{frame}

\begin{frame}{We have the tools!}
\begin{columns}
% Column 1
\begin{column}{0.5\textwidth}
% \center \includegraphics[width=\textwidth]{Illustrations/run0_RBM_color_full_30000.pdf} 
\end{column}
\begin{column}{0.5\textwidth}
\begin{overprint}
\onslide<1>
\begin{itemize}
    \item Graph databases allow us to collect and analyze lineages
    \item Processing and \emph{visualizing} these lineages can help us understand run dynamics
\end{itemize}
\onslide<2>
These visualizations help:
\begin{itemize}
	\item Turn complex runs into ``readable'' ancestry trees
    \item Capture and display key data in nodes and edges
    \item Find valuable properties and patterns in our runs
\end{itemize}
\end{overprint}
\end{column}
\end{columns}
\end{frame}

\begin{frame}{Other tools also exist}
We're not alone:
\begin{itemize}
	\item Bogdan Burlacu has done very similar work with HeuristicLab~\cite{Burlacu:2013:GECCOcomp:new, Burlacu:CIEES:2015}
    \item Tools like ELICIT (Cruz, et al)~\cite{cruz2015elicit} and EpochX (Vaseux, et al)~\cite{vaseux2013event} also support lineage exploration
\end{itemize}

~

All of these are functionality provided by a particular EC system

~

We use:
\begin{itemize}
	\item ``Off-the-shelf'' tools (graph DBs \& visualization libraries)
    \item Could be applied to results from most any EC system
\end{itemize}
\end{frame}

\subsection*{Outline}

\begin{frame}{Outline}
  \tableofcontents[hideallsubsections]
\end{frame}

%~~~~~~~~~~Background~~~~~~~~~~~~~~~~~~~~~~~~~~~~~~~~~~~~~~~~~~~~~~~~~~~~~~~~~~~~~~~~~~~~~~~~~~~~~~~~~~~~~~~

\section{Setup}

\subsection{Run Configuration}

\begin{frame}{GP System}
GP Tools:
\vspace{.5em}
\begin{itemize}
    \setlength\itemsep{.5em}
	\item Clojush implementation of the PushGP system
    % An individual has a linear plush genome
	\item Plush (Linear Push) genomes~\cite{Helmuth:2016:GPTP}
    % Individuals are chosen to reproduce through Lexicase selection.
	\item Lexicase selection~\cite{Helmuth:2014:ieeeTEC}
\end{itemize}
\vspace{1em}
GP Runs:
\begin{itemize}
    \item 1,000 individuals per generation
    \item Up to 300 generations here, 
    	\\ \hspace{.5cm} Over 1,000 generations in other cases
        \\ \hspace{.5cm} 100Ks or millions of individuals
\end{itemize}
\end{frame}

\begin{frame}{Test problem}
Replace-Space-With-Newline software synthesis problem:
\vspace{.5em}
\begin{itemize}
    \setlength\itemsep{.5em}
	\item Two tasks to solve:
    % Use default numbers not in balls
    \setbeamertemplate{enumerate items}[default]
    \begin{enumerate}
    	\item Print input string with all spaces replaced by newlines
        \item Return number of non-space characters in input
    \end{enumerate}
	\item Results in two error values per test case
	\item Used 100 test cases, \\ \hspace{.25cm} So a total of 200 error values stored in vector
\end{itemize}
\vspace{1em}
%\centering
We have hundreds of runs of this problem which we'll use as examples throughout the talk.
\end{frame}

\subsection{Database Configuration}
\begin{frame}{Database schema}
Used TitanDB graph database to store our data.
\vspace{.5cm}
\begin{columns}[t]
% Column 1
\begin{column}{0.45\textwidth}
Nodes = individuals
	\begin{itemize}
	    % Fill the screen more by adding space
    	\setlength\itemsep{.5em}
		\item UUID
        % These aren't stored in the Individual nodes in the DB -- they're used to generate the edges
		% \item UUID of parent(s)
		\item Plush genome
		\item Error vector
		\item Number of children
		\item Number of selections
	\end{itemize}
\end{column}
% Column 2
\begin{column}{0.5\textwidth}
Edges = parent-child relationships
	\begin{itemize}
	    % Fill the screen more by adding space
    	\setlength\itemsep{.5em}
		\item Genetic operator
		\item Damerau-Levenshtein distance
	\end{itemize}
%    \vspace{2cm} % Pushes up second column
\end{column}
\end{columns}
\end{frame}


%~~~~~~~~~~Graph Structure~~~~~~~~~~~~~~~~~~~~~~~~~~~~~~~~~~~~~~~~~~~~~~~~~~~~~~~~~~~~~~~~~~~~~~~~~~~~~~~~~~~~~~
\section{Graph Structure \& Visualization}

\subsection{Graphviz and DOT}

\begin{frame}{Graphviz and DOT for visualization}
\begin{columns}
\begin{column}{.6\textwidth}
\begin{itemize}
    \setlength\itemsep{.5em}
	\item Load run data into database for processing \& analysis
    \item Queries extract relevant data
    \item Scripts generate DOT input
    \item DOT generates PDF output
\end{itemize}
\vspace{0.5cm}

PDF output supports arbitrary zooming \\ \hspace{.5cm} -- huge amounts of detail in paper
\end{column}
\begin{column}{.4\textwidth}
% \includegraphics[width=.9\textwidth]{Illustrations/DOT_Graphiz.pdf}
\end{column}
\end{columns}
\end{frame}

\subsection[Edges]{Edges}

\begin{frame}{Genetic Operators}
Edges represent parent-child relationships
\\
\vspace{.75em}
Styled to indicate genetic operators used to create the child
\\
\vspace{.75em}

Two classes of operators: Alternation \& Mutation
\\
\vspace{.75em}

Four different operators:
\begin{itemize}
% Fill the screen more by adding space
% \setlength\itemsep{.5em} 
% Emphasize displaying variety not details of operators
\item Alternation % A form of Crossover
\item Uniform-Mutation % A form of Point Mutation
\item Alternation \& Uniform-Mutation
\item Uniform-Close-Mutation % Mutates where code blocks end
\end{itemize} 
\end{frame}


\begin{frame}{Edge Coloring \& Style}
	\begin{columns}
		\begin{column}{.55\textwidth}
		\begin{overprint}
			\onslide<1>
				\hspace{.5cm} \textit{Edge Coloring}:  \vspace{.5em}
                \setbeamertemplate{itemize subitem}{}% Remove subitem bullets for this frame
				\begin{itemize}
                % Fill the screen more by adding space
				\setlength\itemsep{1em}
					\item Alternation is 
                    \begin{itemize}
	                    \item \underline{black} and solid
                    \end{itemize}
					\item Uniform-Mutation is
                    \begin{itemize}
	                    \item \underline{\color{orange} orange} and solid
                    \end{itemize}
					\item Alternation \& Uniform-Mutation
                    \begin{itemize}
	                    \item \underline{black} and dashed
                    \end{itemize}
					\item Uniform-Close-Mutation is
                    \begin{itemize}
	                    \item \underline{\color{orange} orange} and dashed
                    \end{itemize}
				\end{itemize} 
			\onslide<2>
				\hspace{.5cm} \textit{Edge Opacity}: Based on \\ \vspace{.00cm}
                \hspace{0.75cm} number of children of \\ \vspace{.00cm}
                \hspace{0.75cm}  the child. \vspace{.5em}
				\begin{itemize}
                % Fill the screen more by adding space
				\setlength\itemsep{1em}
					\item Transparent: Few Children
					\item Opaque: Many Children
				\end{itemize} 
			\onslide<3>
				\hspace{.5cm} Edge Thickness: Based on \\ \vspace{.00cm}
                \hspace{0.75cm} Damerau-Levenshtein distance, \\ \vspace{.00cm}
                \hspace{0.75cm} (measure of genome difference) \vspace{.5em}
				\begin{itemize}
                % Fill the screen more by adding space
				\setlength\itemsep{1em}
					\item Thin: Very Different Genome \\ \hspace{.5cm} -- High DL-Distance 
					\item Wide: Very Similar Genome \\ \hspace{.5cm} -- Low DL-Distance
				\end{itemize}
		\end{overprint}
		\end{column}
		
		\begin{column}{.5\textwidth}
		\begin{overprint}
			\onslide<1-3>
			 	% \includegraphics[width=\textwidth]{Illustrations/family_edges.pdf} 
		\end{overprint}
		\end{column}
\end{columns}
\end{frame}


\begin{frame}{Edge Coloring \& Style Example}
				% \includegraphics[width=\textwidth]{Illustrations/run0_edges_detail.pdf}
\end{frame}


\subsection[Nodes]{Nodes}

\begin{frame}{Node Basics}
Nodes represent an individual or program:
\begin{itemize}
\item Number of selections: Width
\item Number of children: Height
\end{itemize}
\bigskip
\hspace{6.25cm} $\leftarrow$ Number of Selections $\rightarrow$
\begin{tabular}{rc}
\begin{tabular}{r}
\hspace{1cm} $\uparrow$ Number of Children $\downarrow$
\end{tabular}
\begin{tabular}{c}
% \includegraphics[width=.45\textwidth]{Illustrations/node_and_edge.pdf}
\end{tabular}
\end{tabular}
\end{frame}

\begin{frame}{Node Coloring}
% \center
{There are two different techniques we use:  \\ Error Based \& RBM}
% \includegraphics[width=\textwidth]{Illustrations/run0_dual_and_RBM_full.pdf}
\end{frame}

\begin{frame}{Error Based Node Coloring}
	\begin{overprint}
		\onslide<1-2> 			\centering
% \includegraphics[width=.85\textwidth]{Illustrations/run0_bi_color_shaded_percent0.pdf}
	\end{overprint}
	\begin{overprint}
		\onslide<1>
		\centering
		Two colors based on the percent of solved test cases for each task
		\onslide<2>
		\centering
		Additional shading (i.e., darker) for exceptionally high total errors
	\end{overprint}
\end{frame}

\begin{frame}{Restricted Boltzmann Machine Node Coloring}
\centering
%\includegraphics[width=.85\textwidth]{Illustrations/run0_RBM_color_full_30000.pdf} \\
Dimensionality Reduction Problem: 200 test cases $\rightarrow$ 24 bit RGB color
\end{frame}

\begin{frame}{Restricted Boltzmann Machine Node Coloring}
\centering
% \includegraphics[width=.9\textwidth]{Illustrations/run0_RBM_color_full_30000.pdf}
\end{frame}

% \begin{frame}{RBM Example with Run 1: 129 Generations}
% \centering
% \includegraphics[width=.35\textwidth]{Illustrations/run1_RBM_color_full_60000.pdf} 
% \end{frame}

%~~~~~~~~~~What did I learn?~~~~~~~~~~~~~~~~~~~~~~~~~~~~~~~~~~~~~~~~~~~~~~~~~~~~~~~~~~~~~~~~~~~~~~~~~~~~~~~~~~~~~~~

\section{What did we learn?}

\subsection{Hyperselection is common -- is it important?}
\begin{frame}{Hyperselected nodes are highly visible}
% \center \includegraphics[width=\textwidth]{Illustrations/cropped.pdf} \\
In our graphs, we can easily see hyperselection as very wide nodes. 
\end{frame}

%\begin{frame}{Hyperselection is very common}
%\begin{center}
%% \includegraphics[width=0.8 \textwidth]{Illustrations/run0_RBM_color_full_30000.pdf} \\
%Lots of hyperselection -- how important is it to the success of runs?~\cite{Helmuth:2016:GECCO}
%\end{center}
%\end{frame}

\subsection{Filtering ancestry trees}

\begin{frame}{Filtering ancestry trees}
\begin{columns}
% Column 1
\begin{column}{.6\textwidth}
\begin{overprint}
\onslide<1>
We only graph ancestors of ``winning'' individuals.

~

But still a lot of data!
\begin{itemize}
\item 22,435 individuals
\item 35,403 edges
\item (Avg. branching factor of 1.5)
\end{itemize}

\onslide<2>

So we filter out parents who did not contribute much genetic information.

~

Much more ``readable'':
\begin{itemize}
	\item 1,597 individuals (vs 22,435)
    \item 1,795 edges (vs 35,403)
    \item Avg. branching factor of 1.1
\end{itemize}

\onslide<3>
\begin{itemize}
\item Filtering in paper is a bit of a hack
\item Currently (last few weeks) working on tracking every gene (instruction) as it passes from parent to child
\item Allows us to visualize \emph{exactly} the parents that contributed genetic material
\end{itemize}
\end{overprint}

%\begin{itemize}
%\item Numbers on edges indicate similarity.
%\item The smaller the number is, the more similar they are.
%\end{itemize}
\end{column}
% Column 2
\begin{column}{.4\textwidth}
\begin{overprint}
\onslide<1>
% \includegraphics[height=.85\textheight]{Illustrations/run1_RBM_color_full_60000.pdf} 
\onslide<2-3>
% \includegraphics[height=.85\textheight]{Illustrations/run1_RBM_color_filtered_and_full_60000_2.pdf} 
\end{overprint}
\end{column}
\end{columns}
\end{frame}

\begin{frame}{Close-up of filtering}
% \includegraphics[width=\textwidth]{Illustrations/run1_RBM_filtered_and_full_crop.pdf} 
\end{frame}
 
\begin{frame}{Another filtering example}
\center
394 Individuals \hspace{4.5cm} 54 Individuals \\
% \includegraphics[height=.8\textheight]{Illustrations/run0_RBM_color_both_runs_40.pdf} 
\end{frame}

\subsection{Long (nearly) linear chains}
\begin{frame}{Long (nearly) linear chains}
% \includegraphics[width=\textwidth]{Illustrations/run1_RBM_color_filtered_crop.pdf}
% \\
Filtering reveals long nearly linear sequences of pseudo-mutations
\end{frame}

\subsection{Ability to compare runs}

\begin{frame}{Ability to compare runs}
\begin{center}
% \includegraphics[width=0.9\textwidth]{Illustrations/runs_1_99_6_RBM_color_filtered3_rotated.pdf}
\end{center}
\end{frame}

\section{Conclusions}
\begin{frame}{Conclusions}
These visualizations have:
\begin{itemize}
\item Helped turn complex runs into ``readable'' ancestry trees
\item Modified aspects of both nodes and edges to capture key information from the run dynamics
\item Explored ways to filter ancestry trees
\item Highlighted prevalence of hyperselection in lexicase selection runs
\item Highlighted long chains of pseudo-mutations
\item Suggested numerous research directions
\end{itemize}
\end{frame}

\begin{frame}{Future work}
\begin{itemize}
	\item Improve the filtering
    \begin{itemize}
    	\item Currently working to track at the level of genomes (instructions)
        \item Allows more precise tracking of genetic contributions
    \end{itemize}
	\item Enable more dynamic visualizations 
    \begin{itemize}
    	\item Deep Tree exhibit~\cite{block2012deeptree}
        \item Burlacu's work~\cite{Burlacu:CIEES:2015}
    \end{itemize}
    \item Enable more ``on-line'' visualization during runs
\end{itemize}
\end{frame}

\begin{frame}{Thanks!}
\center \Large
Thank you for your time \& attention! \\ \medskip
\includegraphics[width=.1\textwidth]{Illustrations/smile.png} \\ \medskip
\normalsize
~
Special thanks to Bill Tozier, and the folks at the Computational Intelligence Lab at Hampshire College.

~

This material is based upon work supported by the National Science Foundation under 
Grants No. 1129139 and 1331283.

\end{frame}

%~~~~~~~~~~References~~~~~~~~~~~~~~~~~~~~~~~~~~~~~~~~~~~~~~~~~~~~~~~~~~~~~~~~~~~~~~~~~~~~~~~~~~~~~~~~~~~~~~~

\section*{References}

\begin{frame}[allowframebreaks]
\frametitle{References}
\bibliographystyle{abbrv}
\bibliography{GECCO-2016.bib}
\end{frame}

\end{document}

\begin{frame} 
	\frametitle{References} 
	
	\begin{thebibliography}{lskdjf}
	\small 
	\bibitem{Burlacu:2013:GECCOcomp:new}
B.~Burlacu, M.~Affenzeller, M.~	Kommenda, S.~Winkler, and G.~Kronberger.
\newblock Visualization of genetic lineages and inheritance
information in genetic programming.
\newblock In Christian Blum, \emph{et al}, editors, {\em GECCO '13}, pages 1351--1358, New York, NY, USA, 2013.
	
	\bibitem{hyper:2016}
	T.~Helmuth, N.~F.~McPhee, and L.~Spector.
\newblock The Impact of Hyperselection on Lexicase Selection.
\newblock {\em GECCO '16}, Denver, CO, USA 2016.
  
	\bibitem{vis:2016}
	N.~F.~McPhee, M.~Casale, M.~Finzel, T.~Helmuth, and L.~Spector.
\newblock Visualizing genetic programming ancestries using graph databases.
\newblock {\em GECCO '16}, Denver, CO, USA 2016.
	
	\bibitem{gptp:2016}
	N.~F.~McPhee, M.~Finzel, M.~Casale, T.~Helmuth, and L.~Spector.
\newblock A detailed analysis of a PushGP run.
\newblock (Forthcoming) {\em GPTP '16}, Ann Arbor, MI, USA 2016.
	

  
  	\end{thebibliography}
	
\end{frame} 

